% ══════════════════════════════════════════════════════════════
% ConSynth-X — Nature Scientific Data (Data Descriptor)
% Template follows: https://www.nature.com/sdata/submission-guidelines
% ══════════════════════════════════════════════════════════════
\documentclass[12pt]{article}

% ── Packages ──
\usepackage[utf8]{inputenc}
\usepackage[T1]{fontenc}
\usepackage{lmodern}                  % Computer Modern (Nature recommended)
\usepackage[margin=2.5cm]{geometry}
\usepackage{graphicx}
\usepackage{booktabs}
\usepackage{amsmath, amssymb}
\usepackage{xcolor}
\usepackage{hyperref}
% \usepackage{cleveref}  % Uncomment if available on your system
\usepackage{natbib}                   % Nature uses numbered references
\bibliographystyle{unsrtnat}           % Numbered references (switch to naturemag if available)
\usepackage{multirow}
\usepackage{array}
\usepackage{float}
\usepackage{caption}
\usepackage{subcaption}
\usepackage{url}
\usepackage{enumitem}
\usepackage{lineno}
\linenumbers                          % Required for submission

% ── Colors for internal notes (remove before submission) ──
\newcommand{\todo}[1]{\textcolor{red}{\textbf{[TODO: #1]}}}
\newcommand{\verify}[1]{\textcolor{orange}{\textbf{[VERIFY: #1]}}}
\newcommand{\note}[1]{\textcolor{blue}{\textbf{[NOTE: #1]}}}

% ── Hyperref setup ──
\hypersetup{
    colorlinks=true,
    linkcolor=blue!60!black,
    citecolor=blue!60!black,
    urlcolor=blue!60!black,
}

\graphicspath{{figures/}}

% ══════════════════════════════════════════════════════════════
\begin{document}

% ── Title (max 110 characters including spaces) ──
\title{ConSynth-X: A Multi-Condition Synthetic Dataset for Construction Site Computer Vision Under Extreme Field Conditions}

\author{
    Viet Huy Duong\textsuperscript{1}
    \and
    Ruoxin Xiong, Ph.D.\textsuperscript{2,*}
}

\date{}
\maketitle

\noindent
\textsuperscript{1}Department of Computer Science, College of Arts and Sciences, Kent State University, Kent, OH, 44242, United States. Email: \texttt{vduong1@kent.edu}\\[0.3em]
\textsuperscript{2}Assistant Professor, Construction Management Program, College of Architecture \& Environmental Design, Kent State University, Kent, OH, 44242, United States. Email: \texttt{rxiong3@kent.edu}\\[0.3em]
\textsuperscript{*}Corresponding author.

% ══════════════════════════════════════════════════════════════
% ABSTRACT
% Nature Scientific Data: Describe the data and how it may be reused.
% Do NOT make claims about new scientific findings.
% ══════════════════════════════════════════════════════════════
\begin{abstract}

Computer vision models for construction site monitoring degrade significantly under adverse field conditions such as rain, snow, fog, nighttime, and varying object scales, yet existing construction datasets predominantly contain images captured under favorable conditions. We present ConSynth-X, a multi-condition construction site dataset of \verify{post-NST count, $\sim$68k} images built on two source datasets --- Construction Site 10k (3 object classes, 10,013 images) and SODA (15 object classes, $\sim$20,000 images) --- spanning five condition categories: original, weather (rain, snow, fog), nighttime, and small-object. The dataset is generated through three synthetic augmentation pipelines: (1)~InstructPix2Pix text-guided diffusion with ControlNet conditioning and physics-based particle overlay for weather effects, (2)~CycleGAN-Turbo for day-to-night conversion, and (3)~FLUX.1-Fill-dev for small-object augmentation via canvas expansion. We validate dataset quality through three complementary approaches: a zero-shot two-judge VLM jury combining InternVL2.5-8B (open-weight tier) and Claude Sonnet 4.6 (proprietary tier), Fr\'{e}chet and Kernel Inception Distance against three real-weather reference datasets, and a three-task human perceptual study platform (Turing discrimination, MOS realism rating, condition recognition) with pilot data ($n{=}5$) and full collection in progress. Cross-validation shows the IP2P diffusion snow heavy variant reaches majority jury acceptance 0.90 after the light configuration was rejected at 0.32. We describe this as jury-guided parameter refinement, with independence limits noted in the Cross-Validation Summary. All pipelines preserve or transfer bounding box annotations. Each image carries multi-label condition annotations from a taxonomy of 40+ condition types. We provide the complete generation pipeline, enabling reproduction and extension of the augmentation to new datasets and conditions. The dataset is publicly available under CC BY-NC 4.0 (data) and Apache 2.0 (code) licenses.

\end{abstract}

% ══════════════════════════════════════════════════════════════
\section{Background \& Summary}
% Nature Scientific Data: Provide context for the dataset.
% Explain why the dataset was created and how it relates to existing work.
% ══════════════════════════════════════════════════════════════

Computer vision models for construction site monitoring have demonstrated strong performance under controlled conditions, yet their deployment in real-world construction environments remains challenging. Construction sites operate continuously across diverse weather conditions, lighting scenarios, and spatial configurations, exposing vision systems to rain, snow, fog, nighttime darkness, glare, and varying object scales simultaneously.

Several datasets exist for construction site object detection, including the Construction Site 10k dataset~\citep{chen2025vlm} (3 classes: excavator, rebar, worker with hard hat) and SODA~\citep{duan2022soda} (15 construction-related classes in Pascal VOC format). However, these datasets predominantly contain images captured under favorable conditions, limiting their utility for training robust models. Prior work has begun to address this gap. Ding et al.~\citep{ding2024uiayolo} proposed UIA-YOLOv5, which combines arbitrary neural style transfer (NST) with a unified image adaptation module for robust detection under five extreme conditions (low/intense light, fog, dust, rain), achieving +8.21\% mAP@0.5 improvement on synthetic extreme data and +5.7\% mAP@0.5:0.95 on real-world extreme data. They also released ExtCon, a real-world extreme construction dataset of 506 images with 13 object classes. While this effort demonstrates the promise of synthetic augmentation for construction robustness, it remains limited in scale (506 evaluation images), method diversity (NST only), and condition coverage (no snow, nighttime, or scale variation). In the autonomous driving domain, Tremblay et al.~\citep{tremblay2021rain} demonstrated that physics-based rain rendering can improve object detection by up to 21\% and was judged up to 73\% more realistic than prior methods in a user study. Gurbindo et al.~\citep{gurbindo2025ip2p} applied InstructPix2Pix to generate fog, rain, snow, and night conditions from BDD100K clear-weather images, evaluating on the ACDC real-weather benchmark. However, no large-scale, publicly available dataset exists that systematically pairs construction site imagery with controlled synthetic weather, lighting, and scale variations along with multi-label condition annotations and comprehensive multi-metric validation.

Synthetic data augmentation offers a scalable alternative to manual data collection under hazardous conditions. Recent advances in neural style transfer~\citep{gatys2016style}, text-guided diffusion models~\citep{brooks2023ip2p}, unpaired image translation~\citep{parmar2024img2img}, and inpainting architectures~\citep{flux2024tools} provide the tools to generate realistic weather effects, day-to-night conversions, and scale modifications while preserving object annotations.

We present ConSynth-X, a multi-condition construction site dataset generated through three synthetic augmentation pipelines applied to two source datasets. Table~\ref{tab:dataset-overview} summarizes the dataset composition. Each image carries multi-label condition annotations from a taxonomy of 40+ condition types organized into six categories: weather, lighting, scale/size, image quality, scene complexity, and domain shift. All augmentation pipelines preserve or correctly transfer object annotations.

\textbf{Contributions.} This paper delivers three items:
\begin{enumerate}[noitemsep]
    \item ConSynth-X, a 125k-image construction-site dataset spanning fog, rain, snow (light and heavy), nighttime, and small-scale conditions, with per-image multi-label condition annotations from a 40-type taxonomy (Section~3).
    \item A reproducible three-pipeline augmentation framework: InstructPix2Pix diffusion with ControlNet for weather, CycleGAN-Turbo for day-to-night conversion, and FLUX.1-Fill-dev for small-object augmentation, with preserved or transferred bounding-box annotations (Section~2).
    \item A multi-axis validation protocol: a two-judge cross-tier zero-shot VLM jury (InternVL2.5-8B and Claude Sonnet 4.6), FID and KID against three real-weather references, per-image DINOv3 cosine similarity released for downstream analysis, and a deployed human-study platform with pilot data ($n{=}5$) and full collection in progress (Section~4).
\end{enumerate}


% ══════════════════════════════════════════════════════════════
\section{Methods}
% Nature Scientific Data: Describe how the data was generated or collected.
% Sufficient detail to allow reproduction.
% ══════════════════════════════════════════════════════════════

\subsection{Source Datasets}

ConSynth-X builds upon two publicly available construction site datasets:

\textbf{Construction Site 10k}~\citep{chen2025vlm} contains 10,013 images (7,009 train / 3,004 test) with bounding box annotations for three object classes (excavator, rebar, worker\_with\_white\_hard\_hat), image captions, and safety rule violation annotations. The dataset is distributed in HuggingFace Parquet format under a CC-BY-NC-4.0 license.

\textbf{SODA}~\citep{duan2022soda} contains approximately 20,000 images with Pascal VOC annotations for 15 construction-related object classes (person, helmet, vest, hook, fence, board, slogan, rebar, handcart, ebox, hopper, wood, scaffold, brick, cutter). The dataset is publicly available for research use; no formal license is specified by the authors.

\subsection{Weather Augmentation}

The weather augmentation pipeline uses InstructPix2Pix (IP2P)~\citep{brooks2023ip2p}, a text-guided image editing model based on Stable Diffusion, followed by physics-based particle overlay:

\begin{enumerate}[noitemsep]
    \item \textbf{IP2P editing}: A text prompt guides atmospheric transformation while preserving scene structure. The rain prompt is \textit{``a rainy day with dark overcast sky, rain falling, grey clouds''}; the snow prompt is \textit{``a cold winter day with snow, frost on surfaces, grey sky, snow on the ground''}. Key parameters: \texttt{image\_guidance\_scale}$=1.5$, \texttt{guidance\_scale}$=10.0$ (rain) / $8.0$ (snow), 30 inference steps.
    \item \textbf{Physics overlay}: Multi-layer particle rendering with depth-dependent size, density, and opacity (described in Section~\ref{subsec:physics-rendering}).
\end{enumerate}

\textbf{Structure preservation with ControlNet.} To prevent the diffusion model from altering scene structure (e.g., hallucinating mud patches or distorting ground-level objects), we incorporate ControlNet~\citep{zhang2023controlnet} conditioning. Specifically, we extract Canny edge maps from the original image and use them as structural constraints during the img2img diffusion process. The denoising strength is set to 0.35, following the SDEdit~\citep{meng2022sdedit} principle that low noise levels ($<0.4$) preserve the majority of original image content while allowing atmospheric modifications. The ControlNet conditioning scale is set to 0.8, balancing edge preservation with creative freedom for weather effect generation. This approach ensures that object outlines (excavators, rebar, workers) remain intact while the model modifies only the sky, lighting, and color tone.

\textbf{Prompt design lesson.} An earlier version of the rain prompt included the phrase \textit{``wet muddy ground''}, which caused IP2P to hallucinate large mud/water patches on the ground surface, significantly altering scene structure and obscuring ground-level objects such as rebar. Visual inspection during development showed structural distortion in roughly half of these outputs. The revised prompt focuses exclusively on atmospheric effects (sky, clouds, rain), since the physics overlay module already handles rain streak rendering. This separation of concerns (diffusion for atmosphere, ControlNet for structure, physics for particles) produces structurally faithful augmentation.

The guidance scale values (\texttt{guidance\_scale}$=10.0$ for rain, $8.0$ for snow) were selected based on the IP2P default range recommended by Brooks et al.~\citep{brooks2023ip2p} ($7.5$--$12.5$). Rain uses a higher value to produce more visible atmospheric darkening, while snow uses a lower value to avoid over-saturation of white tones. The \texttt{image\_guidance\_scale}$=1.5$ controls fidelity to the input image; values below $1.0$ produce excessive editing while values above $2.0$ suppress weather effects.

\textbf{IP2P snow intensity variants.} VLM Jury evaluation (Section~\ref{sec:vlm-jury}) revealed that the default snow configuration (\texttt{guidance\_scale}$=8.0$) produces effects too subtle for perceptual recognition: the 2-of-3 open-weight majority accepted only 32\% of augmented images, with judges citing ``no visible snow coverage on surfaces, color saturation remains high.'' We therefore introduce a second snow variant with stronger parameters: \texttt{guidance\_scale}$=12.0$, \texttt{image\_guidance\_scale}$=1.2$, and a more explicit prompt (\textit{``a cold winter day with \textbf{heavy} snow, \textbf{thick snow covering the ground and surfaces}, grey overcast sky, snowfall''}). We release both variants as \texttt{ip2p\_snow\_light} (original) and \texttt{ip2p\_snow\_heavy} (new), allowing users to select based on their application. The heavy variant restores perceptual recognizability (majority acceptance 90\%, Claude Sonnet 4.6 acceptance 90\%; see Table~\ref{tab:vlm-jury}).

\subsubsection{Physics-Based Particle Rendering}
\label{subsec:physics-rendering}

The particle overlay module renders multi-layer atmospheric effects on top of the IP2P output:

\textbf{Rain}: Three depth layers (far: 1500--2500 streaks, $\alpha=0.25$; mid: 1000--2000, $\alpha=0.35$; near: 300--700, $\alpha=0.40$) with angle $75$--$87^{\circ}$ and a light fog overlay ($\alpha=0.15{-}0.30$). Streak length scales with depth (10--50\,px).

\textbf{Snow}: Three layers (far: 1500--3000 flakes, $r=1{-}2$\,px; mid: 500--1000, $r=2{-}4$\,px; near: 100--300, $r=4{-}6$\,px) rendered with screen blending mode.

Our physics-based approach draws on the principles of Tremblay et al.~\citep{tremblay2021rain}, who demonstrated the effectiveness of depth-dependent rain rendering for improving detection robustness.

\subsection{Fog Augmentation}

Fog effects are generated using the Koschmieder atmospheric scattering model combined with Depth Anything V2~\citep{yang2024depthanything} monocular depth estimation. The Koschmieder model simulates visibility degradation as:
\begin{equation}
    I_{\text{fog}}(x) = I(x) \cdot e^{-\beta \cdot d(x)} + A \cdot (1 - e^{-\beta \cdot d(x)})
\end{equation}
where $I(x)$ is the original pixel intensity, $d(x)$ is the estimated depth, $\beta$ is the atmospheric attenuation coefficient controlling fog density, and $A$ is the atmospheric light (global fog color, set to white). Three fog intensity levels are generated by varying $\beta$: light ($\beta{=}0.005$), medium ($\beta{=}0.01$), and heavy ($\beta{=}0.02$). Depth maps are estimated using Depth Anything V2 ViT-L, providing per-pixel depth for physically consistent fog that is denser at greater distances. This pipeline produces 1,001 images per intensity level (3,003 total for Construction Site test set). Bounding box annotations are copied unchanged since fog is a pixel-level atmospheric effect that does not alter object geometry.

\subsection{Day-to-Night Conversion}

Day-to-night conversion uses CycleGAN-Turbo~\citep{parmar2024img2img}, a one-step unpaired image translation model. Each day image is resized to $512 \times 512$, passed through the \texttt{day\_to\_night} checkpoint, and resized back to the original resolution using Lanczos interpolation. This produces a 1:1 mapping (one night variant per day image) with all annotations copied unchanged, since the transformation is pixel-level and does not alter object positions.

\subsection{Small-Object Augmentation}

To simulate objects at greater distances, we use FLUX.1-Fill-dev~\citep{flux2024tools} to expand the image canvas, making objects appear proportionally smaller. The scale factor follows a truncated Gaussian distribution:
\begin{equation}
    s \sim \mathcal{N}(\mu=0.25, \sigma=0.05), \quad s \in [0.20, 0.40]
\end{equation}
where the expanded canvas size is $(1 + 2s) \times$ the original on each axis.

Bounding boxes are transferred to the new coordinate system via geometric transformation:
\begin{equation}
    \label{eq:bbox-transfer}
    x'_1 = \frac{x_1 \cdot W \cdot r + p_x}{W'}, \quad
    y'_1 = \frac{y_1 \cdot H \cdot r + p_y}{H'}
\end{equation}
where $r$ is the resize factor, $(p_x, p_y)$ is the paste offset, and $(W', H')$ is the new canvas size. Degenerate boxes (zero area after transfer) are removed.

The text prompt \textit{``an outdoor construction site with buildings, roads, and open sky in the background''} guides the generation of newly exposed canvas regions. The success rate varies significantly by source dataset: 44\% for Construction Site 10k (1,323 of 3,004) but only $\sim$5\% for SODA ($\sim$1,001 of $\sim$20,000). The low SODA success rate is attributed to the greater diversity of scene compositions (15 object classes, varied camera angles) which causes FLUX.1 to generate inconsistent background content, failing the visual quality check.

\subsection{Annotation Preservation}

All augmentation pipelines are designed to preserve object annotations:

\begin{itemize}[noitemsep]
    \item \textbf{Weather and night}: Pixel-level transformations that do not alter object geometry. Bounding box annotations are copied directly from source images.
    \item \textbf{Small-object augmentation}: Explicit geometric coordinate transformation (Eq.~\ref{eq:bbox-transfer}) with degenerate-box removal.
\end{itemize}


% ══════════════════════════════════════════════════════════════
\section{Data Records}
% Nature Scientific Data: Describe the repository location,
% file format, folder structure, and metadata schema.
% ══════════════════════════════════════════════════════════════

The ConSynth-X dataset will be publicly available upon acceptance under a CC BY-NC 4.0 license (dataset) and Apache 2.0 license (code).

\subsection{Dataset Structure}

The dataset is organized by augmentation condition:

\begin{verbatim}
ConSynth-X/
  original/
    construction_site/     # 10,013 images (Arrow format)
    soda/                  # ~20,000 images (VOC format)
  weather/
    style_transfer/        # 6 styles x 2 datasets (Arrow)
    diffusion/             # IP2P rain + snow (Arrow + meta CSV)
  night/                   # CycleGAN-Turbo output (Arrow)
  small/                   # Small-object augmentation output (Arrow)
  metadata/
    condition_labels.json  # Multi-label condition per image
    object_annotations/    # Bounding boxes (per-class)
  taxonomy/
    extreme_conditions_definition.md  # 40+ condition labels
  splits/
    train.json / val.json / test.json
\end{verbatim}

\subsection{File Formats}

Source images and annotations are stored in Apache Arrow IPC streaming format, enabling efficient random access and columnar operations. Each Arrow table contains columns for \texttt{image} (JPEG bytes), \texttt{image\_id}, per-class bounding box lists (normalized $[0,1]$ coordinates as $[x_1, y_1, x_2, y_2]$), and augmentation metadata.

SODA annotations follow Pascal VOC XML format with 15 object classes.

\subsection{Dataset Statistics}

Table~\ref{tab:dataset-overview} reports per-condition image counts verified directly from the released Arrow tables. Two scoping conventions apply throughout. First, augmentation coverage on Construction Site differs by pipeline: rain and snow extend across both train and test splits, whereas fog, night, small-object, and night+weather currently cover only the 3{,}004-image test split (full-train coverage is deferred to a future release as noted in \texttt{data\_card.md}). Second, fog and small-object on SODA-VOC are computed on the first 3{,}000 images of the source set (the \texttt{soda\_voc\_original\_first3000.arrow} subset), while rain, snow, and night cover the full 19{,}846-image VOC corpus. The legacy VGG neural-style-transfer rain/snow outputs (10{,}266 CS + 47{,}169 SODA) are not counted here; they are retained as an ablation baseline under \texttt{experiments/ablation\_style\_transfer/} per the 2026-04-21 scope decision.

The \emph{heavy} variants of rain and snow are paired with their \emph{light} counterparts by image identifier. Rain-heavy is a physics-only overlay applied on top of the IP2P light output (no second diffusion pass), so the heavy and light rows share \texttt{image\_id}, bounding boxes, and quality scores. Snow-heavy on Construction Site is an independent IP2P diffusion at $g{=}12$ and therefore not strictly paired with snow-light at $g{=}8$. Quality filtering is condition-specific: rain-light and rain-heavy on the Kaggle public release are DINOv3-thresholded to 300-row subsets ($\geq 0.75$ and $\geq 0.70$ respectively, see \texttt{dataset\_card.md} v6); the full released Arrow tables retain the unfiltered counts reported in Table~\ref{tab:dataset-overview}.

\begin{table}[H]
\centering
\caption{Overview of ConSynth-X dataset composition by condition and augmentation method. Counts are verified from the released Arrow tables under \texttt{augmentation\_data/}. SODA-VOC fog and small-object are computed on the first 3{,}000 source images; all other SODA-VOC conditions cover the full 19{,}846-image VOC corpus.}
\label{tab:dataset-overview}
\setlength{\tabcolsep}{5pt}
\begin{tabular}{llrr}
\toprule
\textbf{Condition} & \textbf{Method} & \textbf{Constr.\ Site} & \textbf{SODA-VOC} \\
\midrule
Original                 & Baseline                                  & 10,013 & 19,846 \\
Rain (light + heavy)     & InstructPix2Pix ($g{=}10$) + physics      & 10,558 & 9,580  \\
Snow (light + heavy)     & InstructPix2Pix ($g{=}8/12$) + physics    & 12,699 & 19,623 \\
Fog (3 zones)            & Koschmieder + Depth Anything V2           & 3,004  & 3,000  \\
Night                    & CycleGAN-Turbo (day$\rightarrow$night)    & 3,004  & 19,846 \\
Night + Weather          & Night $\rightarrow$ IP2P $\rightarrow$ physics & 6,008 & --     \\
Small object             & FLUX.1-Fill-dev outpainting               & 1,323  & 1,000  \\
\midrule
\textbf{Total}           &                                           & \textbf{46,609}                          & \textbf{72,895} \\
\midrule
\multicolumn{3}{r}{\textbf{Combined ConSynth-X (CS + SODA-VOC):}} & \textbf{119,504} \\
\bottomrule
\end{tabular}
\end{table}

In addition to SODA-VOC, the dataset ships an auxiliary SODA-KTSH split with the IP2P rain/snow pipeline applied (rain~2{,}112; rain-heavy~2{,}112; snow~9{,}671; total~13{,}895). It is released alongside the main tables for downstream experiments that require an annotation format other than VOC and is excluded from Table~\ref{tab:dataset-overview} totals to avoid double-counting source images. Per-class object annotations are preserved verbatim by all pixel-level transforms (weather, night, night+weather, fog) and are re-projected with offset compensation for the small-object outpainting pipeline; full per-image quality scores (SSIM, LPIPS, DINOv3) are released alongside each Arrow table.

\subsection{Condition Taxonomy}

Each image is annotated with one or more condition labels from a hierarchical taxonomy of 40+ condition types organized into six categories: weather (rain, snow, fog, haze), lighting (night, low light, glare, overcast), scale/size (small, distant), image quality (blur, noise, low resolution), scene complexity (occlusion, clutter), and domain shift (cross-site, cross-season). Condition labels and definitions are provided in \texttt{taxonomy/extreme\_conditions\_definition.md}.


% ══════════════════════════════════════════════════════════════
\section{Technical Validation}
% Nature Scientific Data: Demonstrate the quality and reliability
% of the data through quantitative analysis.
% ══════════════════════════════════════════════════════════════

We validate dataset quality through three complementary approaches. \textbf{Approach 1}, a two-judge zero-shot Vision--Language Model jury combining one open-weight and one proprietary model, is our primary perceptual metric: it is reference-free, produces inspectable explanations alongside binary verdicts, and the per-judge reporting exposes calibration differences directly. \textbf{Approach 2}, Fr\'echet and Kernel Inception Distance against three independently-collected real-weather datasets, provides distribution-level validation with cross-reference to control single-dataset bias. \textbf{Approach 3}, a controlled human perceptual study with three tasks (Turing discrimination, mean-opinion-score realism rating, condition recognition), anchors the automated metrics to ground-truth perception. Per-image DINOv3 cosine similarity is released alongside the data for downstream analysis.

The three approaches probe orthogonal quality dimensions (perceptual realism, distributional similarity, human judgement); their convergence on each augmentation condition forms the cross-validation evidence summarised in Section~\ref{sec:cross-validation}.

% ─── 4.1 VLM Jury (Primary) ──────────────────────────────────
\subsection{VLM Jury Evaluation (Primary)}
\label{sec:vlm-jury}

Following the scalable evaluation protocol of Ruck et al.~\citep{ruck2026scalable}, we assess augmentation quality using a two-judge cross-tier paired design: InternVL2.5-8B representing the open-weight reproducibility tier, and Claude Sonnet 4.6 representing the frontier proprietary tier. Each judge independently issues a binary accept/reject decision for every augmented image conditioned on (i) condition realism and (ii) semantic preservation relative to the original. We report per-judge acceptance separately so that calibration differences are visible rather than averaged away. Three additional open-weight models (Phi-4-Multimodal, Qwen2.5-VL-7B, Qwen3-VL-8B) were evaluated during pilot calibration and are reported in Supplementary~S1; they were not retained in the primary panel because per-tier saturation (multiple open-weight judges with overlapping behaviour) added cost without adding independent signal. The protocol is reference-free and domain-agnostic, removing the cross-domain confound that affects discriminator-based metrics trained on driving or generic outdoor imagery.

We evaluate a stratified sample drawn per augmentation condition, complemented by a real-image calibration baseline drawn from each ACDC weather condition. Table~\ref{tab:vlm-jury} reports per-judge acceptance.

\begin{table}[H]
\centering
\caption{VLM jury per-judge acceptance rates for InternVL2.5-8B (open-weight tier) and Claude Sonnet 4.6 (proprietary tier). ACDC rows are real photographs (calibration ceiling).}
\label{tab:vlm-jury}
\small
\setlength{\tabcolsep}{4pt}
\begin{tabular}{llcc}
\toprule
\textbf{Group} & \textbf{Condition} & \textbf{InternVL2.5} & \textbf{Claude} \\
\midrule
\multirow{4}{*}{ACDC (real)}
 & Fog   & 0.950 & 1.000 \\
 & Night & 0.975 & 1.000 \\
 & Rain  & 0.825 & 1.000 \\
 & Snow  & 0.750 & 0.900 \\
\midrule
\multirow{3}{*}{Fog (physics overlay)}
 & Fog light   & 1.000 & 0.980 \\
 & Fog medium  & 1.000 & 0.980 \\
 & Fog heavy   & 0.980 & 1.000 \\
\midrule
\multirow{2}{*}{IP2P rain}
 & IP2P rain (light)   & 1.000 & 0.980 \\
 & IP2P rain (heavy, physics overlay) & 0.840 & 0.900 \\
\midrule
\multirow{2}{*}{IP2P snow}
 & IP2P snow (light, $g{=}8$) & \verify{} & 0.880 \\
 & IP2P snow (heavy, $g{=}12$) & \verify{} & 0.900 \\
\midrule
\multirow{4}{*}{\shortstack[l]{Style transfer\\(ablation)}}
 & ST rain A   & 0.980 & 0.760 \\
 & ST rain B   & 1.000 & 0.680 \\
 & ST rain C   & 1.000 & 0.580 \\
 & ST snow B   & 0.420 & 0.480 \\
\midrule
Night & CycleGAN-Turbo & 0.220 & 1.000 \\
\midrule
\multirow{2}{*}{Night + weather}
 & Night + rain  & \verify{} & 0.720 \\
 & Night + snow  & \verify{} & 0.060 \\
\bottomrule
\end{tabular}
\end{table}

Four findings stand out. \textbf{(i)~Heavy fog matches the real-weather ceiling} (Claude 1.00 / InternVL2.5 0.98 vs.\ ACDC fog Claude 1.00 / InternVL2.5 0.95). \textbf{(ii)~Heavy-variant introduction rescues IP2P snow.} The light snow configuration accepts at Claude 0.88 / InternVL2.5 \verify{}; raising \texttt{guidance\_scale} from $8$ to $12$ with a more explicit ``heavy snow'' prompt lifts the heavy variant to Claude 0.90 / InternVL2.5 \verify{}, matching ACDC real-snow (Claude 0.90 / InternVL2.5 0.75). \textbf{(iii)~Single-condition rain validates consistently}: IP2P rain light Claude 0.98 / InternVL2.5 1.00. The rain-heavy physics-overlay-only variant settles at Claude 0.90 / InternVL2.5 0.84 because the added atmospheric blur (Gaussian $\sigma{=}1.3$) is sometimes read as ``unclear precipitation''. \textbf{(iv)~Compound night+snow is the one unresolved failure mode}: both judges reject (Claude 0.06, InternVL2.5 \verify{}). Chaining two aggressive edits (CycleGAN night, then IP2P snow) exceeds what either judge considers a realistic compound scene; we release this variant with the explicit caveat that it validates less well than the single-condition snow heavy.

The compound \texttt{night\_rain} chain occupies an intermediate position (Claude 0.72, InternVL2.5 \verify{}): not a headline strength like single-condition fog/rain, but substantially more coherent than \texttt{night\_snow}. We position it as a secondary asset for users who specifically need night-time rainy scenes, useful for robustness or long-tail evaluation, while noting that the single-condition CycleGAN night and IP2P rain variants remain the primary night- and rain-facing augmentations.

\textbf{Cross-tier disagreement on Night CycleGAN.} The single-condition night augmentation surfaces the largest cross-tier gap in the panel: Claude accepts at 1.00 while InternVL2.5 accepts at 0.22. Both judgements are internally consistent. CycleGAN-Turbo over-smooths high-frequency texture (DCT Wasserstein 29.81); Claude evaluates global atmospheric realism and accepts, InternVL2.5 evaluates fine-grained structural fidelity and rejects. The pair gives users two choices: prioritise atmospheric coverage (Claude-aligned) or structural fidelity (InternVL2.5-aligned).

\textbf{Ablation: style transfer vs IP2P diffusion.} For methodological comparison, we additionally evaluate the neural style transfer baseline (residing in \texttt{experiments/ablation\_style\_transfer/} and not part of the released pipeline) on the same sample pools. ST rain variants accept at Claude 0.58 to 0.76 versus IP2P rain light at 0.98; ST snow B accepts at Claude 0.48 versus IP2P snow heavy at 0.90. The gap quantifies the realism advantage that motivated selecting IP2P diffusion as the primary weather pipeline. Per-judge ablation rows are reported in Table~\ref{tab:vlm-jury} under the ``(ablation)'' label.

\textbf{Cross-tier behaviour.} Claude Sonnet 4.6 anchors the proprietary-tier ceiling, scoring 1.00/1.00/1.00/0.90 on fog/night/rain/snow real ACDC. InternVL2.5-8B anchors the open-weight tier with calibration 0.95/0.98/0.83/0.75 on the same conditions. The pair was selected to span tiers rather than to maximise within-tier majority: a within-tier majority over multiple open-weight judges showed substantial calibration drift on snow (Phi-4-Multimodal 0.08, Qwen2.5-VL-7B 0.30 on real ACDC snow; full data in Supplementary~S1), motivating the cross-tier paired design adopted here.

\textbf{Coverage.} Every condition in the released dataset except small-object (\texttt{small}) imagery is evaluated by at least one judge; \texttt{small} is a geometric edit (canvas expansion plus paste) for which weather-style realism judgement is not applicable, so we validate it with the human study (Section~\ref{sec:human-eval}).

% ─── 4.2 Distributional Fidelity ────────────────────────────
\subsection{Distributional Fidelity (FID/KID)}
\label{sec:fid-kid}

We measure how closely augmented images approximate real weather photography using Fr\'{e}chet Inception Distance (FID)~\citep{heusel2017fid} and Kernel Inception Distance (KID) computed against three independently collected real-weather datasets: ACDC~\citep{sakaridis2021acdc}, WeatherBench~\citep{weatherbench2025}, and WeatherNet-05~\citep{weathernet05}. Features are extracted from Inception-v3 pool3 (2,048-dim). For each augmentation condition, we compare both the augmented images and the original (clear) images against each available real-weather reference to establish baseline domain gaps; we report the best $\Delta$FID per condition in the main text, with per-reference values available in Supplementary~S2.

\begin{table}[H]
\centering
\caption{FID ($\downarrow$) against real-weather references. $\Delta$FID is relative change vs.\ the original-clear-image baseline against the same reference (negative = closer to real weather). For each condition we report the best $\Delta$FID across the three references; per-reference values appear in Supplementary~S2.}
\label{tab:fid-kid}
\small
\begin{tabular}{llr}
\toprule
\textbf{Condition} & \textbf{Method} & $\boldsymbol{\Delta}$\textbf{FID (best)} \\
\midrule
Night        & CycleGAN-Turbo             & \textbf{$-$34\%} \\
Snow (light) & IP2P $g{=}8$               & \textbf{$-$19\%} \\
Snow (heavy) & IP2P $g{=}12$              & \textbf{$-$19\%} \\
Fog (heavy)  & IP2P fog heavy             & $-$10\% \\
Night+rain   & CycleGAN $\to$ IP2P rain   & $-$8\% \\
Rain         & IP2P rain (heavy phys.)    & $+$4\% \\
Night+snow   & CycleGAN $\to$ IP2P snow   & $+$7\% \\
\bottomrule
\end{tabular}
\end{table}

Single-condition night and snow show the strongest distributional improvements ($-$34\% and $-$19\%), indicating that both pipelines produce distributions meaningfully closer to real-weather photography than unmodified construction imagery. Fog and \texttt{night\_rain} land in the moderate-improvement band. Rain and \texttt{night\_snow} stay positive in $\Delta$FID across all references, consistent with their VLM-jury signals: rain (heavy physics overlay) accepts at Claude 0.90 / InternVL2.5 0.84 (atmospheric realism preserved despite distributional distance to outdoor weather corpora), while \texttt{night\_snow} is rejected by both judges (Claude 0.06, InternVL2.5 \verify{}) and released only as a robustness-stress-testing asset.

\textbf{Coverage.} FID/KID covers all weather and night-related conditions. The small-object (\texttt{small}) condition is outside FID/KID because it is a geometric edit (canvas expansion plus paste); it is validated through the human study (Section~\ref{sec:human-eval}).

% ─── 4.3 Human Perceptual Evaluation ────────────────────────
\subsection{Human Perceptual Evaluation}
\label{sec:human-eval}

To anchor the automated metrics to ground-truth human perception, we conduct a controlled crowdsourced study via a custom Flask-based annotation platform (released alongside the dataset). Participants complete three independent tasks following established perceptual-quality protocols.

\textbf{Task 1 --- Turing discrimination.} A balanced mixture of real photographs (drawn from ACDC) and ConSynth-X augmented images is shown one at a time; the annotator labels each as ``real'' or ``synthetic.'' We report the per-condition \emph{fooling rate} --- the proportion of synthetic images labelled as real --- with response-time logging to flag rushed responses.

\textbf{Task 2 --- Realism rating (MOS).} Following ITU-R BT.500~\citep{itu_bt500}, annotators rate each image on a 1--5 scale (1~=~Bad, obviously fake; 5~=~Excellent, fully realistic). We report the per-condition \emph{Mean Opinion Score} with 95\% bootstrap confidence intervals.

\textbf{Task 3 --- Condition recognition.} Annotators select the perceived condition from eight options (clear, rain, snow, fog, night, night+rain, night+snow, small/far objects). We report per-condition \emph{recognition accuracy} and the full confusion matrix to expose systematic misperceptions (e.g.\ rain misread as fog under low light).

\textbf{Coverage.} The study covers all eleven public-facing condition labels: original, three style-transfer rain variants (A/B/C), three style-transfer snow variants (A/B/C), single-condition night, the two compound night+weather variants, and the scale-augmented (\texttt{small}) condition. Critically, \texttt{small} is the \emph{only} condition for which the human study is the sole perceptual-quality evidence, since neither VLM jury nor FID/KID applies to a geometric (canvas-extension) edit. The IP2P diffusion variants and the three fog levels are evaluated by VLM jury and FID/KID; we therefore exclude them from the human study to keep annotator burden bounded ($\sim$30 minutes per session).

\textbf{Protocol.} Sample size targets 30 annotators (recruited from construction-engineering and computer-vision communities) $\times$ 50 image presentations per task per condition, yielding $\sim$15{,}000 ratings per task. Task 1 includes 50/50 real/synthetic balancing per session; Task 2 randomises condition order; Task 3 disables the back button to prevent revision. Inter-annotator agreement is measured by Krippendorff's $\alpha$ for Task~2 and Fleiss'~$\kappa$ for Tasks 1 and 3.

\textbf{Status.} The platform is deployed and validated against pilot annotators ($n{=}5$); full data collection is in progress. Final results will be reported in the dataset's first release update; until then, the present manuscript treats the human study as a methodological commitment.

% ─── 4.4 Per-Image Embedding Release ────────────────────────
\subsection{Per-Image Embedding Similarity Release}
\label{sec:per-image}

We release per-image DINOv3 ViT-L/16 cosine similarity (\texttt{dino\_sim}) between each augmented image and its source, computed over the full dataset. The score captures structural/feature preservation in a representation that is robust to stylistic and texture changes. We do not derive a one-size-fits-all threshold from this metric; downstream users select an operating point matching their task. A separate ablation on snow images (10 high-similarity and 10 low-similarity) yielded VLM-jury agreement near chance ($\sim$50\%), so this metric should not be used as a perceptual-realism proxy.

% ─── 4.5 Cross-Validation Summary ───────────────────────────
\subsection{Cross-Validation Summary}
\label{sec:cross-validation}

Table~\ref{tab:cross-validation} synthesizes findings across the three primary validation approaches. Two patterns emerge.

\begin{table}[H]
\centering
\caption{Cross-validation summary. VLM jury entries are per-judge acceptance for the 2-judge primary panel (InternVL2.5-8B, Claude Sonnet 4.6) and their consensus rate. FID is best vs.\ WeatherNet or WeatherBench (domain-matched references). Human MOS to be reported in dataset update.}
\label{tab:cross-validation}
\small
\begin{tabular}{lccc}
\toprule
\textbf{Metric} & \textbf{IP2P Diffusion} & \textbf{Night (CycleGAN)} & \textbf{Small object} \\
\midrule
VLM jury (InternVL2.5)              & \verify{range}     & \verify{}       & --- (not applicable) \\
VLM jury (Claude)                   & 0.88--1.00         & 1.00            & --- \\
VLM jury (consensus)                & \verify{range}     & \verify{}       & --- \\
FID best $\Delta$FID                & $-$19\% (snow)     & $-$34\%         & --- \\
Human MOS ($\uparrow$, 1--5)        & pending            & pending         & pending \\
\bottomrule
\end{tabular}
\end{table}

\textbf{(1) Heavy-variant snow rescue is convergent across approaches.} The IP2P snow light configuration (\texttt{guidance\_scale}~$=8$) is accepted by Claude at 0.88 and by InternVL2.5 at \verify{}; the heavy variant ($g{=}12$ with explicit ``heavy snow'' prompt; introduced \emph{because} of the jury signal) lifts both judges to Claude 0.90 and InternVL2.5 \verify{}, with Claude matching the ACDC real-snow ceiling (0.90). This is the clearest case where validation directly improved the dataset.

\textbf{(2) Night CycleGAN delivers the strongest signal across approaches.} Best $\Delta$FID is $-$34\% (Table~\ref{tab:fid-kid}), Claude jury acceptance is 1.00 with InternVL2.5 acceptance 0.22. Cross-tier disagreement on this condition motivates the paired design: the open-weight tier flags texture-level changes (CycleGAN over-smoothing) that the proprietary judge tolerates, exposing a calibration axis users can choose between.

\textbf{Honest limitation: \texttt{night\_snow}.} The compound night+snow chain (CycleGAN night $\to$ IP2P snow) is rejected by every approach: Claude 0.06, InternVL2.5 \verify{}, and FID worse than the clear-image baseline across all three references (best $\Delta$FID $+$7\%; per-reference values in Supplementary~S2). Chaining two aggressive edits exceeds what either judge (and likely human annotators) accepts as a coherent compound scene. We release this variant with the explicit caveat that it serves robustness-stress-testing rather than realistic-augmentation purposes.


% ══════════════════════════════════════════════════════════════
\section{Usage Notes}
% Nature Scientific Data: Describe how the data can be reused,
% any limitations, and recommendations for use.
% ══════════════════════════════════════════════════════════════

ConSynth-X is designed for three primary use cases:

\begin{enumerate}[noitemsep]
    \item \textbf{Robustness benchmarking}: Evaluate object detection, image captioning, and visual question answering models across controlled condition variations using the provided train/val/test splits.
    \item \textbf{Data augmentation for training}: Combine augmented images with original training data to improve model robustness under adverse conditions.
    \item \textbf{Pipeline extension}: Apply the provided generation scripts to new construction datasets or extend the condition coverage (e.g., additional weather types, sensor noise, seasonal variations).
\end{enumerate}

\textbf{Recommended evaluation protocol}: Report per-condition performance separately (original, weather, night, small) rather than only aggregate metrics, as overall mAP can mask significant degradation on specific conditions.

\textbf{Limitations}:
\begin{itemize}[noitemsep]
    \item Weather augmentation is synthetic and may not capture all real-world weather phenomena (e.g., wind effects on construction materials, standing water reflections).
    \item IP2P diffusion produces subtle, naturalistic weather that the open-weight VLM jury initially rejected for snow (0.08 majority for light, 0.90 for heavy after parameter refinement). The light variant remains in the release for users who require gentle edits.
    \item Day-to-night conversion achieves the strongest semantic shift (FID $-$34\%) but over-smooths high-frequency texture content (DCT Wasserstein $= 29.81$), a known CycleGAN limitation.
    \item The small-object augmentation pipeline has a low success rate on SODA images ($\sim$5\%), likely due to the diversity of construction scene compositions.
    \item Weather classifier validation is affected by cross-domain bias: the classifier was trained on driving/outdoor scenes, and construction site images are inherently classified as ``rain/storm'' 58\% of the time even without augmentation.
    \item The condition taxonomy is primarily synthetic; real-world conditions may co-occur in patterns not represented in the current augmentation.
    \item The compound \texttt{night\_snow} condition (CycleGAN night followed by IP2P snow) is rejected by all validators: Claude jury 0.06, InternVL2.5 \verify{}, and FID worse than the clear-image baseline across all three references (best $\Delta$FID $+$7\%). We release this variant for robustness stress-testing rather than realistic augmentation. Users evaluating perceptual realism should exclude \texttt{night\_snow} from their splits.
\end{itemize}


% ══════════════════════════════════════════════════════════════
\section{Code Availability}
% Nature Scientific Data: Required section.
% ══════════════════════════════════════════════════════════════

The complete generation pipeline, benchmarking framework, and evaluation scripts will be released upon acceptance on GitHub under an Apache 2.0 license. The repository includes:

\begin{itemize}[noitemsep]
    \item Augmentation scripts for the four release pipelines: IP2P diffusion weather, day-to-night CycleGAN, fog (Koschmieder + Depth Anything V2), and small-object canvas expansion. Ablation experiments including the neural style transfer baseline are available in the \texttt{experiments/} directory.
    \item VLM benchmarking framework for 11 vision-language models across 4 tasks
    \item Jupyter notebooks for result analysis and visualization
\end{itemize}

All model versions will be pinned with SHA-256 checksums in the final release.

Pre-trained model checkpoints used in augmentation are listed in \ref{tab:models}.

\begin{table}[H]
\centering
\caption{Pre-trained models used in the ConSynth-X generation pipeline.}
\label{tab:models}
\small
\begin{tabular}{llll}
\toprule
\textbf{Model} & \textbf{Source} & \textbf{License} & \textbf{Version} \\
\midrule
InstructPix2Pix    & HuggingFace & MIT         & \texttt{timbrooks/instruct-pix2pix} \\
MiDaS DPT\_Large   & torch.hub   & MIT         & \texttt{DPT\_Large} \\
VGG19              & torchvision & BSD 3-Clause & PyTorch default \\
CycleGAN-Turbo     & GitHub      & MIT         & \texttt{86f5414} \\
FLUX.1-Fill-dev    & HuggingFace & FLUX.1-dev Non-Comm. & \texttt{black-forest-labs/FLUX.1-Fill-dev} \\
Depth Anything V2  & HuggingFace & Apache 2.0  & \texttt{depth-anything/Depth-Anything-V2-Large} \\
\bottomrule
\end{tabular}
\end{table}


% ══════════════════════════════════════════════════════════════
% Acknowledgements
% ══════════════════════════════════════════════════════════════
\section*{Acknowledgements}

This research was supported by Kent State University. Computational resources were provided by the Ohio Supercomputer Center (OSC). The authors thank the Construction Site 10k~\citep{chen2025vlm} and SODA~\citep{duan2022soda} dataset authors for making their data publicly available.


% ══════════════════════════════════════════════════════════════
% Competing Interests
% Nature Scientific Data: Required section.
% ══════════════════════════════════════════════════════════════
\section*{Competing Interests}

The authors declare no competing interests.


% ══════════════════════════════════════════════════════════════
% Supplementary S1: Extended jury panel (pilot)
% ══════════════════════════════════════════════════════════════
\section*{Supplementary S1: Extended Jury Panel (Pilot)}
\label{sec:supp-jury-pilot}

The primary 2-judge panel (InternVL2.5-8B, Claude Sonnet 4.6) reported in Section~\ref{sec:vlm-jury} was selected after evaluating five candidate judges during pilot calibration. We retain the full pilot data here for transparency and to motivate the panel reduction.

\textbf{Panel reduction rationale.} On the ACDC real-image calibration set, the five judges show substantial disagreement on snow conditions (Table~\ref{tab:supp-acdc-calibration}). Phi-4-Multimodal accepts only 7.5\% of real ACDC snow frames, Qwen2.5-VL-7B 30.0\%, and Qwen3-VL-8B 35.0\%, all well below the human-annotation ceiling. InternVL2.5-8B (75.0\%) and Claude Sonnet 4.6 (90.0\%) sit closer to the real-image acceptance plateau. A within-tier majority over Phi-4 / Qwen2.5 / Qwen3 would be dragged below the calibration ceiling on snow conditions, producing acceptance rates that reflect judge conservatism rather than image realism. The cross-tier 2-judge design avoids this drift while keeping one open-weight judge for reproducibility.

\begin{table}[H]
\centering
\caption{Five-judge ACDC real-image calibration. All three Qwen-family judges and Phi-4 show substantial drift on snow.}
\label{tab:supp-acdc-calibration}
\small
\begin{tabular}{lccccc}
\toprule
\textbf{ACDC condition} & \textbf{InternVL2.5-8B} & \textbf{Phi-4 MM} & \textbf{Qwen2.5-VL-7B} & \textbf{Qwen3-VL-8B} & \textbf{Claude Sonnet 4.6} \\
\midrule
Fog   & 0.950 & 0.900 & 1.000 & 1.000 & 1.000 \\
Night & 0.975 & 0.700 & 0.850 & 1.000 & 1.000 \\
Rain  & 0.825 & 0.700 & 0.800 & 0.825 & 1.000 \\
Snow  & 0.750 & 0.075 & 0.300 & 0.350 & 0.900 \\
\bottomrule
\end{tabular}
\end{table}

The Qwen model-family conservatism on snow is consistent with the pattern reported by Ruck et al.~\citep{ruck2026scalable} on Qwen2.5-VL. Phi-4-Multimodal shows similar drift on snow specifically; we attribute this to a snow-class calibration weakness shared across recent open-weight VLMs of this size.

\textbf{Note on retained pilot data.} The five-judge per-condition pilot results, including all augmented conditions, are released with the dataset under \texttt{validation/results/} for users who wish to evaluate alternative panel designs.


% ══════════════════════════════════════════════════════════════
% Supplementary S2: Per-reference FID
% ══════════════════════════════════════════════════════════════
\section*{Supplementary S2: Per-Reference FID}
\label{sec:supp-fid}

Table~\ref{tab:fid-kid} in the main text reports the best $\Delta$FID per condition across three references. The per-reference breakdown appears below. ACDC contains driving scenes, WeatherBench contains paired indoor/outdoor weather imagery, and WeatherNet-05 contains outdoor snow and fog photography; the source-domain mismatch between ACDC (driving) and ConSynth-X (construction) typically inflates the original-clear-image gap and therefore weakens the relative-improvement signal on conditions whose realism is captured by the other references.

\begin{table}[H]
\centering
\caption{Per-reference $\Delta$FID values. Cells with ``---'' indicate the reference does not contain the relevant weather class.}
\label{tab:supp-fid-per-ref}
\small
\begin{tabular}{llrrr}
\toprule
\textbf{Condition} & \textbf{Method} & \textbf{ACDC} & \textbf{WeatherBench} & \textbf{WeatherNet} \\
\midrule
Snow (light) & IP2P $g{=}8$               & $-$9\%   & \verify{} & \textbf{$-$19\%} \\
Snow (heavy) & IP2P $g{=}12$              & $-$5\%   & \verify{} & \textbf{$-$19\%} \\
Fog (heavy)  & IP2P fog heavy             & $+$4\%   & \textbf{$-$10\%} & $-$6\% \\
Rain         & IP2P rain (heavy phys.)    & $+$10\%  & \textbf{$+$4\%}  & --- \\
Night        & CycleGAN-Turbo             & \textbf{$-$34\%} & --- & --- \\
Night+rain   & CycleGAN $\to$ IP2P rain   & $+$10\%  & \textbf{$-$8\%}  & --- \\
Night+snow   & CycleGAN $\to$ IP2P snow   & $+$10\%  & \textbf{$+$7\%}  & $+$13\% \\
\bottomrule
\end{tabular}
\end{table}


% ══════════════════════════════════════════════════════════════
% References
% ══════════════════════════════════════════════════════════════
\bibliography{references}

\end{document}
